\chapter{Conclusion}
\label{chap:conclusion}
\emergencystretch=3em\relax

This work began from a mismatch stated in Chapter~\ref{chap:context}. ENISA's
\emph{Best Practices for Cyber Crisis Management} describes sixteen best
practices distributed across the four phases of a crisis, and the NIS2
directive expects regulated entities to demonstrate a crisis-management
capability rather than simply assert one, yet neither text provides a way to
measure how far an organisation currently stands from that standard. That
chapter turned the mismatch into three constraints that shaped everything built
afterwards: the result has to survive scrutiny from a client's own governance
and security leadership, the referential rather than the tool has to remain the
source of truth, and the instrument has to be usable by a consultant leading a
client conversation rather than only by someone with deep prior expertise in
the ENISA text.

The answer built during this internship starts with the referential itself.
\texttt{referential.\allowbreak yaml} encodes the sixteen best practices as
forty-eight questions across four phases weighted equally at twenty-five per
cent each, a maturity scale from zero to four, evidence caps that bound a
declared level by the strongest evidence the respondent claims, four
foundational gates, consistency penalties between related practices, and an
N/A option that is neutral to the score. Making that file the only
specification of the framework, principle P4 in Chapter~\ref{chap:method}, is
what allows the framework to be reviewed by reading one artefact instead of
reconstructing it from scattered code. On top of it sits a pure Python engine
with a fixed pipeline order, evidence cap, phase-weighted aggregation,
foundational gate, then consistency penalty, described in
Chapter~\ref{chap:scoring-engine}. The engine went through a methodological
audit conducted independently of its own implementation, whose recommendations
were integrated into the code before any platform work began, with the
referential incremented from 1.0.0 to 1.1.0 to record the resulting change of
contract.

Around that core, Chapter~\ref{chap:platform} describes the platform that makes
the engine usable: a FastAPI backend and a Next.js frontend, with a single
translation seam, \texttt{scoring\_\allowbreak bridge.\allowbreak py}, as the
only place allowed to turn a stored assessment into the engine's input, so that
two endpoints can never disagree on the same assessment's score. Real
authentication was built during the internship, with invitation-based account
creation, server-side cookie sessions, three roles, per-owner isolation of
assessments, and access control centralised in one module rather than
re-implemented per route. The local AI stage runs entirely through Ollama, sits
strictly downstream of a fully computed result, receives no organisation name,
identifier or justification text, has its output normalised and checked against
the referential before anything reaches the user, and degrades to the
referential's own static recommendations on any failure. The whole is covered
by 141 automated tests, seventy for the scoring engine and seventy-one for the
API and the backend integration, reported in Chapter~\ref{chap:validation}.

Chapter~\ref{chap:state-of-the-art} situated this against what already exists,
and the space it identified was genuinely empty. General maturity instruments
such as NIST CSF or ISACA's CMMI Cybermaturity Platform treat crisis management
as one dimension among many, sector-specific tools are tied to a single
regulatory context, and academic proposals remain frameworks rather than
deployed instruments; ENISA's own maturity work targets CSIRTs and national
authorities rather than organisation-level crisis preparedness. No published
instrument operationalises the ENISA 2024 framework into a scorable
self-assessment for organisations subject to NIS2. The three positions that
define this project, ENISA operationalisation, a strict separation between
deterministic scoring and AI-generated narrative, and local inference, are the
three the state of the art left converging on an unoccupied point. What this
work contributes is one concrete occupation of that point, built and tested
rather than specified.

That contribution should be read within its limits, which
Chapter~\ref{chap:scoring-engine} and Chapter~\ref{chap:validation} state at
length. The assessment remains declarative, and the evidence caps only bound
how far a favourable answer can inflate a score. The thresholds the model
relies on come from expert judgement rather than from calibration on data. No
inter-rater study was run, and the AI layer was shown to be safe to expose
without being shown to produce good advice. The platform has also not yet been
used on a real client engagement.

What this project gave me, beyond the artefact, is a particular way of
approaching an engineering problem. The value of the scoring engine lies as
much in what it is forbidden to do as in what it computes, and designing a
component around its own constraints, rather than around its features, was new
to me. So was containing a non-deterministic component inside a system that
still has to be defensible, which turned out to be less a question of trusting
the model and more a question of building everything around it on the
assumption that it will occasionally be wrong. The independent audit taught me
that a suite of green tests can be reassuring precisely where it is blind, and
the documentation-first discipline discussed in Chapter~\ref{chap:pm} taught me
its own value largely through its absence at the start of the project, since
the scoring engine, the densest part of the work in methodological decisions,
is the one part whose reasoning I later had to reconstruct rather than consult.

Several directions remain open, and Chapter~\ref{chap:impact} sets them out in
detail rather than here: article-level mapping to NIS2 and ISO/IEC 27035,
retrieval over the source texts, tracking a client's maturity over successive
assessments, and the infrastructure work a real deployment would require. The
one worth closing on is more general. The pattern this project instantiates,
computing a defensible result with conventional logic first and letting a
language model explain or prioritise that result afterwards, owes nothing to
cyber crisis management in particular. It applies wherever an organisation
needs an auditable output alongside natural-language assistance, which is a
large share of what compliance and maturity work actually consists of. Whether
it generalises is a question the next project to try it will answer. This
thesis instantiates that pattern once, in one domain, and demonstrates that
doing so is compatible with keeping every piece of client data inside the
organisation's own infrastructure.